\subsection{The assembled theorem}\hypertarget{tag-01PE}{}\label{fam:section:A.5}

\begin{theorem}\hypertarget{tag-01Q6}{}\label{fam:src:A.5.1}\label{fam:st:A.5.1}\label{wb:pshift:thm:classification}
\StatementRecord{fam:src:A.5.1}{proved}{01Q6}{\SeriesRef{app:rederive:fam-src-A.5.1}; \citet[Theorem~1.4, p.~3]{DuEtAl2026v1}}

Fix the original cusp-regular period solution, the clockwise peripheral marking, the prescribed periodic triangular cusp fan, the affine vectors $(\varepsilon,-\varepsilon')$, and the untwisted cusp collar. For every $c,c'\in\mathcal P$, the constructed compact complex threefolds satisfy
\[
 X_c\simeq X_{c'}\quad\Longleftrightarrow\quad c'-c\in2\mathbb Z.
\]
Every biholomorphism is over the identity of the original $\mathbb P^1$. The proper holomorphic family of their isomorphism classes has parameter $0<|w|<e^{-\pi M}$ and nonzero Kodaira--Spencer class at every point.
\end{theorem}



The standard classification is \citet[Theorem~1.4, p.~3]{DuEtAl2026v1}. A complete re-derivation is given in Appendix~\ref{app:rederive:fam-src-A.5.1}.






The tangent dimensions and Kuranishi germ are re-derived from \citet[Theorem~6.12, pp.~81--82, and Proposition~6.16, p.~85]{DuEtAl2026v1}. A complete re-derivation is given in Appendix~\ref{app:rederive:fam-src-A.5.2}.


\paragraph{Scope and a precise changed-input test.}
The hypotheses of \ref{fam:src:A.5.1}--\ref{fam:src:A.5.2} fix the original period system and
multiplicities $(3,4)$; \ref{fam:section:B.2} treats its arbitrary affine and cusp data,
and \ref{fam:section:B.3} classifies its stipulated higher-multiplicity variants.
A concrete extension beyond these hypotheses begins with
$y^2=x^3+tx+1$ in \cite[Sections 2--3 and Proposition 6.1]{ViaclovskyTwo2026}.
Already one rank-one filling over a specified root $p_1$ of
$4t^3+27=0$ replaces the hexagonal normalization calculation by
\[
 \begin{gathered}
 E'=\mathbb C^*/\langle c_0\rangle,\qquad
 R=\mathbb P_{E'}(\mathcal O_{E'}\oplus L),\\
 R=(\mathbb P^1\times\mathbb C^*)/
       \langle(z,v)\mapsto(\gamma_0z,c_0v)\rangle.
 \end{gathered}
\]
with its boundary sections identified by $v\mapsto\beta_0v$.
Here $c_0,\gamma_0,\beta_0$ are exactly the nonzero central multipliers
in that paper's equation (3.1). The invariant nodal volume form there
trivializes the ambient canonical line, so the branch-normal map has
quotient $\nu_*K_R$ by the same local coefficient map as
\eqref{fam:eq:ab-normal-quotient}. The ruling has
$\pi_*\mathcal O_R=\mathcal O_{E'}$ and $R^1\pi_*\mathcal O_R=0$;
Leray and Serre duality therefore give $h^\bullet(K_R)=(0,1,1)$.
The first unevaluated global datum for transferring the deformation
argument is the connecting map from this single
$H^1(R,K_R)=\mathbb C$ summand into $H^2(T_X(-\log D))$, where
$D$ is the union of the three reduced singular fibres of that model.
Its direct attack is to compute the boundary-translation Cech cocycle
using the full holomorphic unit functions in that equation (3.1),
their central values $(\gamma_0,c_0,\beta_0)$, and the corresponding
vertical frames.
The global identification argument also requires the integral
centralizer and collar defects of that model's four monodromy matrices.
These are specified substitutions to perform for the stated one-fibre
test; no classification or deformation conclusion for that changed
period system is included in \ref{fam:src:A.5.1}--\ref{fam:src:A.5.2}.

